\section{Hot-carrier cascades from first principles}\label{app:E.13}

An electron given \SI{50}{eV} in silicon ends, a picosecond later, as about 13 electron--hole pairs at the band edges and some \SI{35}{eV} of phonons. Which share goes where, and how much the pair number fluctuates from one event to the next, is decided by a competition between two processes at every energy along the way. One is impact ionization, in which the hot carrier gives part of its energy to a valence electron and lifts it across the gap; the other is phonon emission, in which it gives a few tens of meV to the lattice. \Cref{sec:18.1} writes the cascade as a branching process whose kernels are master-formula rates with the carrier as the projectile, and \cref{sec:18.2} defines the mean energy per pair $\varepsilon_{\mathrm{eh}}$ and the Fano factor $F_{\mathrm{eh}}$ as its first two moments. The kernels are computed from the band structure, the screened interaction and the electron--phonon vertex, and silicon is the case where both measurements and a calibrated phenomenological model exist to compare against.

\subsection{Computed quantities}

For a seed class $\Lambda$, the set of initial electron and hole states that the primary excitation produces (\cref{sec:18.3}), the calculation produces the distribution $P_{\mathrm{casc}}(N_{\mathrm{eh}}|\Lambda)$ of the final number of pairs. Its moments are the measured quantities,
\begin{equation}
  \varepsilon_{\mathrm{eh}} = \frac{E_0}{\langle N_{\mathrm{eh}}\rangle}, \qquad
  F_{\mathrm{eh}} = \frac{\langle N_{\mathrm{eh}}^2\rangle - \langle N_{\mathrm{eh}}\rangle^2}{\langle N_{\mathrm{eh}}\rangle},
  \label{eq:E.13-moments}
\end{equation}
for a seed of total energy $E_0$, in the limit of large $E_0$ where both become independent of it. The cascade needs the quasiparticle bands $\varepsilon_{n\vect{k}}$ from the band edges up to the highest energy a carrier reaches, tens of eV for the cascades of \cref{chap:18}; the impact-ionization rate $\varGamma_{\mathrm{ii}}(n\vect{k})$ of a carrier in state $n\vect{k}$, with the distribution of its three products over final states; and the phonon emission and absorption rate $\varGamma_{\mathrm{ph}}(n\vect{k})$, with its final states and phonon branches. The cascade Monte Carlo then follows each carrier from state to state until every carrier is below the impact-ionization threshold $\theta_{\mathrm{ii}}\Eg$ (\cref{sec:18.1}), after which it can only lose energy to phonons and cannot add a pair.

\subsection{Place on the hierarchy}

Each kernel is a golden-rule rate at a definite place on the truncation label $\mathfrak{T}$ of \cref{sec:8.5}. The impact-ionization rate is the decay of a quasiparticle into a lower-energy quasiparticle and an electron--hole pair through the dynamically screened interaction $W^{\mathrm{scr}}$ of \cref{sec:4.8}; with $W^{\mathrm{scr}}$ in the random-phase approximation, its many-body approximation is $\mathcal{X} = \mathrm{RPA}$, applied to quasiparticle rather than Kohn--Sham energies. The phonon rate is one-phonon emission or absorption, $n_g = 2$, with $n_B = 0$. The cascade itself is a sequence of independent golden-rule events: the carrier is on shell between events, and interference between successive events is dropped. This is the same incoherent limit as the transport of \cref{app:E.12}, and it holds when the total scattering rate of a carrier is small compared with its energy above the band edge, the condition that the off-shell parameter be small. A rate of \SI{e14}{s^{-1}} corresponds to a width $\varGamma = \SI{0.066}{eV}$ in natural units, so a carrier \SI{5}{eV} above the conduction-band minimum meets the condition even at rates ten times larger. A carrier a few tens of meV above the impact-ionization threshold, whose energy relative to that threshold is comparable to the width, does not, and the threshold region is where the incoherent picture is least reliable.

\subsection{Quasiparticle bands}

The impact-ionization threshold is set by the gap, and the rate just above it rises steeply with the energy above threshold, so the gap must be right before anything else is. Semilocal density-functional theory places the conduction bands of silicon too low, with an indirect gap of about half the measured \SI{1.17}{eV} (\cref{app:E.1}). Quasiparticle corrections in the $G_0W_0$ approximation (\cref{app:E.3}) repair the gap and the ordering~\cite{hybertsen1986electron}. A self-consistent variant gives \SI{1.23}{eV} for silicon without spin--orbit coupling, and a single empirical scaling of the self-energy correction brings the gap to experiment and the other tabulated conduction-band energies of silicon to within \SI{0.2}{eV} of measurement~\cite[Table~I]{kotani2010impact}. A rigid scissor shift fixes the gap by assuming that the correction is the same for every conduction state, which has to be checked over the tens of eV that the cascade spans.

The band window is large. The number of bands below an energy $E$ above the band bottom is about the number of plane waves in a sphere of radius $k_E = (2m_eE)^{1/2}$ in the unit cell, $\Omega k_E^3/(6\pi^2)$. For silicon, with $\Omega = a^3/4 = \SI{40.05}{\angstrom^3}$, this is about 8 bands up to \SI{20}{eV}, 32 up to \SI{50}{eV} and 90 up to \SI{100}{eV}. All of these bands, on a dense $\vect{k}$ grid, enter both kernels.

\subsection{The impact-ionization kernel}

The rate at which a quasiparticle in an empty state $n\vect{k}$ decays is set by the imaginary part of its self-energy $\Sigma$ (\cref{sec:4.7}). In the $GW$ approximation, $\Sigma = \ii\GF W^{\mathrm{scr}}$, the imaginary part for an unoccupied state is~\cite[Eqs.~(1)--(2)]{kotani2010impact}
\begin{gather}
  \varGamma_{\mathrm{ii}}(n\vect{k}) = 2z_{\mathrm{qp}}\,|\Imag\Sigma_{n\vect{k}}|, \nonumber\\
  \Imag\Sigma_{n\vect{k}} = -\pi\sum_{n'\vect{k}'}\iint\dd^3r\,\dd^3r'\,
  \psi^*_{n\vect{k}}(\vect{r})\psi_{n'\vect{k}'}(\vect{r})\nonumber\\
  \times\Imag W^{\mathrm{scr}}(\vect{r}, \vect{r}', \varepsilon_{n\vect{k}} - \varepsilon_{n'\vect{k}'})\,
  \psi^*_{n'\vect{k}'}(\vect{r}')\psi_{n\vect{k}}(\vect{r}'),
  \label{eq:E.13-imsigma}
\end{gather}
with the sum over empty final states $n'\vect{k}'$ between the Fermi level and $\varepsilon_{n\vect{k}}$, and $z_{\mathrm{qp}}$ the quasiparticle residue of \cref{eq:E.3-qp}, typically about 0.8~\cite[Sect.~II]{kotani2010impact}. The physical content is in $\Imag W^{\mathrm{scr}}$: at energy transfer $\omega = \varepsilon_{n\vect{k}} - \varepsilon_{n'\vect{k}'}$ it is the spectral weight of the excitations the carrier can create, which below the plasmon energy are electron--hole pairs and above it are mostly plasmons. In silicon the plasmon is at \SI{16.6}{eV}~\cite[Table~I]{ramanathan2020ionization}, and for a carrier more than that above the band edge \cref{eq:E.13-imsigma} includes plasmon emission automatically; the plasmon then decays into a pair.

The golden-rule route writes the rate directly as a golden-rule sum over the two electrons and one hole of the final state. It uses the matrix element of the screened interaction between the initial and final pairs, direct and exchange terms combined as $|M|^2 = \tfrac12(|M_{\mathrm{D}}|^2 + |M_{\mathrm{E}}|^2 + |M_{\mathrm{D}} - M_{\mathrm{E}}|^2)$, and momentum conservation $\vect{k} = \vect{k}' + \vect{k}_1 - \vect{k}_2$~\cite{kane1967electron}, \cite[Eq.~(5)]{kotani2010impact}. For weak $\Imag W^{\mathrm{scr}}$ the two expressions differ by the factor $z_{\mathrm{qp}}$ and by the exchange term, which \cref{eq:E.13-imsigma} omits; together these make the two rates differ by a factor between 0.5 and 1~\cite[Sect.~II]{kotani2010impact}. Omitting exchange alone multiplies the rate by no less than 0.75~\cite[Sect.~II]{kotani2010impact}. On the logarithmic scale on which impact-ionization rates are usually compared this is a small difference. For $F_{\mathrm{eh}}$, which depends on how the rate rises above threshold relative to the phonon rate, it is not negligible, and the choice belongs among the stated approximations of a calculation.

The cascade needs more than the total rate. When an ionization event is selected, the energies and momenta of the three products are drawn from the integrand of the rate, resolved by final state. The kernel is then a table of probabilities. For each initial $(n, \vect{k})$ it holds the probabilities of the final $(n', \vect{k}')$ of the scattered carrier and of the transfer $(\vect{q}, \omega)$ to the pair; the pair's electron and hole are then drawn with the weight of the joint density of states at that $(\vect{q}, \omega)$. The screened interaction enters at all transfers up to the carrier energy, so $W^{\mathrm{scr}}$ is needed with full frequency dependence from 0 to tens of eV and with local fields, which is the dielectric calculation of \cref{app:E.6} carried to high frequency.

\subsection{The electron--phonon kernel}

The phonon rate of a carrier in state $n\vect{k}$ is the golden-rule rate of \cref{sec:6.2},
\begin{equation}
  \begin{split}
  \varGamma_{\mathrm{ph}}(n\vect{k}) = 2\pi\sum_{m\nu}\int\frac{\dd^3Q}{V_{\mathrm{BZ}}}\,|g_{mn\nu}(\vect{k}, \vect{Q})|^2
  \Bigl[&(n_{\vect{Q}\nu} + 1)\,\delta(\varepsilon_{n\vect{k}} - \varepsilon_{m,\vect{k}+\vect{Q}} - \omega_{\vect{Q}\nu})\\ &+ n_{\vect{Q}\nu}\,\delta(\varepsilon_{n\vect{k}} - \varepsilon_{m,\vect{k}+\vect{Q}} + \omega_{\vect{Q}\nu})\Bigr],
  \end{split}
  \label{eq:E.13-gph}
\end{equation}
with $V_{\mathrm{BZ}}$ the zone volume, the emission term carrying $n_{\vect{Q}\nu} + 1$ and the absorption term $n_{\vect{Q}\nu}$. At the millikelvin temperatures of a detector, $n_{\vect{Q}\nu} = 0$ for the optical and zone-boundary phonons that dominate, and only emission survives. The vertex $g_{mn\nu}$ comes from density-functional perturbation theory (\cref{app:E.4}) and is interpolated to fine grids with Wannier functions (\cref{app:E.5})~\cite{ponce2016epw,giustino2017electron}.

The catch for a cascade is the energy range. The Wannier set that interpolates the band edges, eight $sp^3$-like orbitals per silicon cell spanning the valence bands and the lowest four conduction bands (\cref{app:E.5}), does not represent states tens of eV up, and its interpolated $g$ is meaningless there. Either the Wannier manifold is enlarged, with disentanglement over a window that contains tens of entangled bands, or the vertex is computed directly on a coarser grid at high energy, where the rates are smoother functions of $\vect{k}$ because the density of final states averages over many bands. Above a few eV the phonon rate is dominated by the emission of optical and intervalley phonons of about \SI{60}{meV}~\cite[Sect.~II\,B]{ramanathan2020ionization}, so the energy lost per phonon event is nearly fixed, and what the calculation must get right is the rate, not the detailed spectrum.

\subsection{Full-band cascade Monte Carlo}

With the two kernels tabulated, the cascade is sampled as in full-band device simulation~\cite{fischetti1988monte}, without a field. Each carrier is a state $(n, \vect{k})$; its total rate is $\varGamma_{\mathrm{tot}} = \varGamma_{\mathrm{ii}} + \varGamma_{\mathrm{ph}}$. The time to its next event is drawn by self-scattering: a constant ceiling $\varGamma_{\mathrm{max}} \ge \varGamma_{\mathrm{tot}}$ is chosen, a time is drawn from $\varGamma_{\mathrm{max}}\ee^{-\varGamma_{\mathrm{max}}t}$, and at that time a real event is accepted with probability $\varGamma_{\mathrm{tot}}/\varGamma_{\mathrm{max}}$, otherwise nothing happens~\cite{jacoboni1983monte}. Without a field the times only order the events and do not change the outcome, but they are kept because the same code with a field gives the drift of \cref{app:E.12}. An accepted event is impact ionization with probability $\varGamma_{\mathrm{ii}}/\varGamma_{\mathrm{tot}}$, which replaces the carrier by three, or phonon emission, which lowers its energy by $\omega_{\vect{Q}\nu}$ and moves it to $m, \vect{k}+\vect{Q}$. Holes are followed in the valence bands with their own kernels, and the hole band width, about \SI{12}{eV} in silicon, limits the energy a hole can take~\cite[Sect.~II\,A]{ramanathan2020ionization}. A carrier that falls below $\theta_{\mathrm{ii}}\Eg$ is removed from the list with its pair counted. The cascade ends when the list is empty.

The seed matters at low energy. \Cref{sec:18.1} explains why the cascade must start from the final states of the primary process, with their momentum and band structure, rather than from a carrier of given energy. The energy sharing between the first electron and hole sets how much energy is above threshold, and below about \SI{20}{eV} in silicon different sharing assumptions give different pair yields~\cite[Sect.~IV\,B]{ramanathan2020ionization}. The seed distribution of \cref{sec:18.3} is drawn into the initial carrier list, and $P_{\mathrm{casc}}(N_{\mathrm{eh}}|\Lambda)$ inherits the seed class $\Lambda$.

\subsection{The single-parameter model}

Silicon has a cascade model that reproduces the measured pair energy, and it is the benchmark. The model of Alig and Bloom, as implemented by Ramanathan and Kurinsky~\cite{ramanathan2020ionization}, replaces both kernels by free-particle forms with a single optical phonon of energy $\omega_{\mathrm{AB}} = \SI{63}{meV}$. For a carrier at energy $E$ above its band edge, the ratio of the phonon and impact-ionization rates is then~\cite[Eq.~(8)]{ramanathan2020ionization}
\begin{equation}
  \frac{\varGamma_{\mathrm{ph}}(E)}{\varGamma_{\mathrm{ii}}(E)} = A_{\mathrm{AB}}\,\frac{105}{2\pi}\,\frac{(E - \omega_{\mathrm{AB}})^{1/2}}{(E - \Eg)^{7/2}} ,
  \label{eq:E.13-ab}
\end{equation}
with energies in eV and one constant $A_{\mathrm{AB}}$, the ratio of the squared phonon and Coulomb matrix elements times phase-space factors, and the probability that the next event is an ionization is $[1 + \varGamma_{\mathrm{ph}}/\varGamma_{\mathrm{ii}}]^{-1}$~\cite[Eq.~(10)]{ramanathan2020ionization}. The value $A_{\mathrm{AB}} = 5.2$, in the units of the source (quoted there as eV$^2$, with which the ratio in \cref{eq:E.13-ab} carries a leftover eV$^{-1}$), was chosen in the original model to reproduce the measured pair energy at high energy, and hot-electron injection data confirm the model up to \SI{5}{eV}~\cite[Sect.~II and Table~1]{ramanathan2020ionization}. It makes the ionization probability per event 0.05 at \SI{3}{eV}, 0.36 at \SI{5}{eV}, 0.77 at \SI{8}{eV} and 0.88 at \SI{10}{eV}, from \cref{eq:E.13-ab} with $\Eg = \SI{1.17}{eV}$. Phonon emission wins below about \SI{6}{eV}, where the two probabilities cross, and the many \SI{63}{meV} losses a carrier suffers on its way down from there are what puts $\varepsilon_{\mathrm{eh}}$ at three times the gap.

Simulated with this model, the high-energy moments follow the linear fits $\varepsilon_{\mathrm{eh}} = 1.7\Eg + 0.084A_{\mathrm{AB}} + 1.3$ and $F_{\mathrm{eh}} = -0.028\Eg + 0.0015A_{\mathrm{AB}} + 0.14$, with energies in eV and $A_{\mathrm{AB}}$ in the same units~\cite[Eqs.~(12)--(13)]{ramanathan2020ionization}. At \SI{123}{K}, where the gap is \SI{1.160}{eV}~\cite[Eq.~(11)]{ramanathan2020ionization}, they give $\varepsilon_{\mathrm{eh}} = \SI{3.71}{eV}$ and $F_{\mathrm{eh}} = 0.115$. The measurement at the same temperature with a charge-resolving silicon sensor is $\varepsilon_{\mathrm{eh}} = \SI{3.749 \pm 0.001}{eV}$ and $F_{\mathrm{eh}} = 0.119 \pm 0.002$~\cite{rodrigues2021absolute}. The pair energy agrees to \SI{1}{\percent} because $A_{\mathrm{AB}}$ was tuned to it; the Fano factor is a prediction and falls two standard deviations low. This is the comparison the first-principles cascade must improve on. Since its kernels contain no fitted constant, matching $\varepsilon_{\mathrm{eh}}$ is a test of the kernels, and $F_{\mathrm{eh}}$ is a second, independent one.

\subsection{Convergence parameters}

For the bands: the quasiparticle correction, the number of bands (\cref{eq:E.13-moments} needs every band a carrier can reach) and the $\vect{k}$ grid. For $W^{\mathrm{scr}}$: the plane-wave cutoff of the dielectric matrix, which must resolve momentum transfers up to about twice the largest carrier momentum, the number of bands in the polarizability, the frequency grid from 0 to the largest carrier energy, and the $\vect{q}$ grid. For the kernels: the $\vect{k}$ grid on which $\varGamma_{\mathrm{ii}}$ and $\varGamma_{\mathrm{ph}}$ are tabulated, and the broadening of the energy-conserving delta functions. Near threshold the impact-ionization rate is a sum over few final states that conserve both energy and momentum. On a finite grid it scatters strongly from one $\vect{k}$ to the next; with 1728 $\vect{k}$ points in the zone the numerical error is about a factor two at rates below \SI{e10}{s^{-1}}~\cite[Fig.~2]{kotani2010impact}. The scatter is a physical feature of the band structure smeared by the grid, and it averages out at higher energies, where many final states contribute~\cite[Sect.~III]{kotani2010impact}.

The statistical parameter is the number of simulated cascades $N_{\mathrm{casc}}$. The mean $\langle N_{\mathrm{eh}}\rangle$ converges quickly, since its relative error is $[F_{\mathrm{eh}}/(\langle N_{\mathrm{eh}}\rangle N_{\mathrm{casc}})]^{1/2}$. The variance converges slowly: for a nearly Gaussian pair number, the relative error of a sample variance is $(2/N_{\mathrm{casc}})^{1/2}$. Resolving $F_{\mathrm{eh}}$ to $\pm0.002$ at 0.119, the precision of the silicon measurement, needs a relative error of \SI{1.7}{\percent} and therefore $N_{\mathrm{casc}} \approx 2/0.017^2 \approx 7\times10^3$ cascades at each seed energy, and more if the comparison is to be made at the two-standard-deviation level.

\subsection{Cost}

The dominant cost is the full-frequency screened interaction. For carriers up to \SI{50}{eV} in silicon, the momentum transfers reach $2k_E = \SI{7.2}{\angstrom^{-1}}$, and a dielectric matrix that resolves them holds $N_G \approx \Omega(2k_E)^3/(6\pi^2) \approx 260$ plane waves. Stored at 200 frequencies on the 29 irreducible points of an $8^3$ $\vect{q}$ grid, in double-precision complex numbers, the matrices $W^{\mathrm{scr}}_{\vect{G}\vect{G}'}(\vect{q}, \omega)$ occupy $29\times260^2\times200\times\SI{16}{byte} \approx \SI{6}{GB}$, and computing them requires the polarizability summed over about 30 empty bands on a $\vect{k}$ grid of similar density. For carriers up to \SI{100}{eV} the band count triples, $N_G$ grows to about 730 and $N_G^2$ eightfold. The kernel tables add a sum over final states for every tabulated initial state: with $N_k$ points and $N_b$ bands it scales as $N_k^2N_b^2$ before symmetry, which for the 1728-point grid above and 30 bands is about $3\times10^9$ terms, each a contraction with $W^{\mathrm{scr}}$. The phonon kernel is a standard Wannier-interpolated calculation (\cref{app:E.5}), cheaper than the electronic one once the high-energy vertex problem is solved. The cascade Monte Carlo is cheap: a \SI{5}{keV} seed makes some 1300 pairs and a few times $10^4$ phonon events, so that $10^4$ cascades amount to some $6\times10^8$ events, negligible next to the tabulation of the kernels.

\subsection{Pitfalls}

The gap must match the detector temperature, \SI{1.17}{eV} at millikelvin against \SI{1.12}{eV} at room temperature~\cite[Table~I]{ramanathan2020ionization}, since every rate near threshold depends on the energy above $\theta_{\mathrm{ii}}\Eg$. Two double countings are easy to make: the $GW$ rate already contains plasmon emission, so no plasmon channel is added; and the first-principles phonon rate of \cref{eq:E.13-gph} replaces, and is not added to, a phenomenological one. Two models that reproduce the same $\varepsilon_{\mathrm{eh}}$ do not validate each other's $F_{\mathrm{eh}}$, and the agreement of the cascade of \cref{sec:18.4} with the single-parameter model is agreement of two models, not independent support.

\subsection{Silicon}

\Cref{tab:E-cascade-si} summarizes the silicon inputs, the reference values and the sizes derived in this section.

\begin{table}[t]
  \caption{The hot-carrier cascade in silicon. Band and matrix sizes are free-electron estimates from the text; the model values follow from \cref{eq:E.13-ab} and its fits at \SI{123}{K}~\cite{ramanathan2020ionization}; the measurement is at \SI{123}{K} and \SI{5.9}{keV}~\cite{rodrigues2021absolute}.}
  \label{tab:E-cascade-si}
  \centering
  \begin{tabular}{p{0.43\linewidth}p{0.45\linewidth}}
    \toprule
    Quantity & Silicon \\
    \midrule
    Gap at low temperature; quasiparticle gap without spin--orbit & \SI{1.17}{eV}; \SI{1.23}{eV}~\cite{kotani2010impact} \\
    Plasmon energy; hole band width & \SI{16.6}{eV}; \SI{12}{eV}~\cite{ramanathan2020ionization} \\
    Bands up to 20, 50, \SI{100}{eV} & 8, 32, 90 \\
    $N_G$ and storage of $W^{\mathrm{scr}}$ for \SI{50}{eV} carriers & 260; \SI{6}{GB} (29 $\vect{q}$, 200 frequencies) \\
    $\vect{k}$ grid; threshold-region error & 1728 points; factor 2 below \SI{e10}{s^{-1}}~\cite{kotani2010impact} \\
    Ionization probability per event at 3, 5, \SI{8}{eV} & 0.05, 0.36, 0.77 (\cref{eq:E.13-ab}) \\
    $\varepsilon_{\mathrm{eh}}$, $F_{\mathrm{eh}}$: model & \SI{3.71}{eV}, 0.115 \\
    $\varepsilon_{\mathrm{eh}}$, $F_{\mathrm{eh}}$: measured & \SI{3.749 \pm 0.001}{eV}, $0.119 \pm 0.002$ \\
    Cascades for $\delta F_{\mathrm{eh}} = 0.002$ & $7\times10^3$ per seed energy \\
    \bottomrule
  \end{tabular}
\end{table}
